% figure 
\begin{figure}[!htb] 
	\centering
		\resizebox{\textwidth}{!}{
			\subcaptionbox{Estimation satisfaisante (cycle 16 de la \cref{fig_unp_position_test_trajectory}). A la fin de la fenêtre temporelle de \SI{200}{\milli\second}, l'erreur est submillimétrique (\SI{0.5}{\milli\meter}) pour une amplitude de mouvement supérieure à \SI{8}{\centi\meter}, ce qui correspond à une erreur relative inférieure à \SI{1}{\percent}.\label{subfig_fig_test_burdet_2000_success}}
			{
				\tikzsetnextfilename{fig_test_burdet_trajectory_prediction_a}
				\begin{tikzpicture}
				\centering
				%
				
				\begin{axis}[%
					name=ax1,
					width=.45\linewidth,
					height=.3\linewidth,
					scale only axis,   
					xlabel={temps (\SI{}{\milli\second})},
					ylabel={Position (\SI{}{\centi\meter})},
					axis x line = bottom,
					axis y line = left, 
					]
					% first 15 candidates
					\foreach \i in {2,3,...,16} {			
						\addplot [
							color=blue,
							solid,
							opacity=0.3,
							]
						table [col sep=comma, x=idx, y index=\i] {misc/test1_burdet/candidates_subsamp.csv};
					}
					%
					\foreach \i in {50,51,...,65} {			
						\addplot [
						color=blue,
						smooth,
						opacity=0.3,
						]
						table [col sep=comma, x=idx, y index = \i] {misc/test1_burdet/candidates_subsamp.csv};
					}
					%
					\foreach \i in {150,151,...,165} {			
					\addplot [
						color=blue,
						smooth,
						opacity=0.3,
						]
						table [col sep=comma, x=idx, y index = \i] {misc/test1_burdet/candidates_subsamp.csv};
					}
					\addlegendentry{Candidats};
					%
					\addplot [
						color=black,
						smooth,
						line width=1.pt
						]
					table [col sep=comma, x expr=\coordindex, y=actual_trajectory] {misc/test1_burdet/best_candidates.csv};
					%
					\addplot [
						color=red,
						dashed,
						line width=1.pt
						]
					table [col sep=comma, x expr=\coordindex, y=burdet_candidate] {misc/test1_burdet/best_candidates.csv};
					
					\coordinate (c1) at (axis cs:580, \pgfkeysvalueof{/pgfplots/ymin});
					\coordinate (c2) at (axis cs:580, \pgfkeysvalueof{/pgfplots/ymax});				
					\coordinate (d1) at (axis cs:480, \pgfkeysvalueof{/pgfplots/ymin});
					\coordinate (d2) at (axis cs:480, \pgfkeysvalueof{/pgfplots/ymax});			
					\coordinate (e1) at (axis cs:780, \pgfkeysvalueof{/pgfplots/ymin});
					\coordinate (e2) at (axis cs:780, \pgfkeysvalueof{/pgfplots/ymax});
	
					\draw [dashed] (c1) -- (c2);
					\draw [dashed] (d1) -- (d2);
					\draw [dashed] (e1) -- (e2);
		
					\draw[>=latex, <->] ($(d1) + (rel axis cs:0,0.45)$) -- node [midway,above,align=center,fill=white,yshift=.35cm] {\footnotesize apprentissage\\ \footnotesize\SI{100}{\milli\second}} ($(c1) + (rel axis cs:0,0.45)$);
					\draw[>=latex, <->] ($(c1) + (rel axis cs:0,0.45)$) -- node [midway,above,align=center,yshift=.15cm] {\footnotesize \SI{200}{\milli\second}} ($(e1) + (rel axis cs:0,0.45)$);
					
					%define coordinates of rectangle
					\coordinate (r1) at (axis cs:600,0.5);
					\coordinate (r2) at (axis cs:850,1.5);
					\coordinate (r3) at (axis cs:600,1.5);
					\coordinate (r4) at (axis cs:850,0.5);			
					\draw (r1) rectangle (r2);
					
				\end{axis}
				%
				\begin{axis}[%
					name=ax2,
					width=.5\linewidth,
					height=.4\linewidth,
					scaled ticks=false,   
					xlabel={temps (\SI{}{\milli\second})},
					axis x line = bottom,
					axis y line = left, 
					xmin=600,xmax=850,
					at={($(ax1.south east)+(1cm,0)$)},
					]
					%
					\addplot [
						color=black,
						solid,
						line width=1.pt,
						]
					table [domain=600:850, col sep=comma, x expr=\coordindex, y=actual_trajectory] {misc/test1_burdet/best_candidates.csv};
					\addlegendentry{Mesures};
	%				%
					\addplot [
						color=red,
						dashed,
						line width=1.pt,
						]
					table [domain=600:850, col sep=comma, x expr=\coordindex, y=burdet_candidate] {misc/test1_burdet/best_candidates.csv};
					\addlegendentry{Meilleur cand.};		
					\coordinate (e1) at (axis cs:780, \pgfkeysvalueof{/pgfplots/ymin});
					\coordinate (e2) at (axis cs:780, \pgfkeysvalueof{/pgfplots/ymax});
					
					\draw [dashed] (c1) -- (c2);
					\draw [dashed] (d1) -- (d2);
					\draw [dashed] (e1) -- (e2);
					
				\end{axis}
				% draw dashed lines from rectangle in first axis to corners of second
				\draw [dashed] (r4) -- (ax2.south west);
				\draw [dashed] (r2) -- (ax2.north west);
				
				\end{tikzpicture}
			}
		}
	\\
		\resizebox{\textwidth}{!}{
			\subcaptionbox{Estimation peu satisfaisante (cycle 25 de la \cref{fig_unp_position_test_trajectory}). A la fin de la fenêtre temporelle de \SI{200}{\milli\second}, l'erreur est de \SI{5.9}{\milli\meter} pour une amplitude de mouvement supérieure à \SI{8}{\centi\meter}, ce qui correspond à une erreur relative d'environ \SI{7}{\percent}.\label{subfig_fig_test_burdet_2000_failure}}
			{
				\tikzsetnextfilename{fig_test_burdet_trajectory_prediction_b}
				\begin{tikzpicture}
				\centering
				%
				%\LoadCSV {mycsv} {misc/test1_burdet/best_candidates.csv}
				%\indexP{\GetCSV{mycsv}{2}{3}} 
				
				\begin{axis}[%
				name=ax1,
				width=.45\linewidth,
				height=.3\linewidth,
				scale only axis, 
				xlabel={temps (\SI{}{\milli\second})}				    
				ymin=-0.05,
				ylabel={Position (\SI{}{\centi\meter})},
				axis x line = bottom,
				axis y line = left, 
				]
				% first 15 candidates
				\foreach \i in {2,3,...,16} {			
					\addplot [
					color=blue,
					smooth,
					opacity=0.3,
					]
					table [col sep=comma, x=idx, y index=\i] {misc/test2_burdet/candidates_subsamp.csv};
				}
				%
				\foreach \i in {50,51,...,65} {			
					\addplot [
					color=blue,
					smooth,
					opacity=0.3,
					]
					table [col sep=comma, x=idx, y index = \i] {misc/test2_burdet/candidates_subsamp.csv};
				}
				%
				\foreach \i in {150,151,...,165} {			
					\addplot [
					color=blue,
					smooth,
					opacity=0.3,
					]
					table [col sep=comma, x=idx, y index = \i] {misc/test2_burdet/candidates_subsamp.csv};
				}
				%
				\addplot [
				color=black,
				solid,
				line width=1.pt
				]
				table [col sep=comma, x expr=\coordindex, y=actual_trajectory] {misc/test2_burdet/best_candidates.csv};
				%
				\addplot [
				color=red,
				dashed,
				line width=1.pt
				]
				table [col sep=comma, x expr=\coordindex, y=burdet_candidate] {misc/test2_burdet/best_candidates.csv};
				%				
				\coordinate (c1) at (axis cs:438, \pgfkeysvalueof{/pgfplots/ymin});
				\coordinate (c2) at (axis cs:438, \pgfkeysvalueof{/pgfplots/ymax});				
				\coordinate (d1) at (axis cs:338, \pgfkeysvalueof{/pgfplots/ymin});
				\coordinate (d2) at (axis cs:338, \pgfkeysvalueof{/pgfplots/ymax});			
				\coordinate (e1) at (axis cs:638, \pgfkeysvalueof{/pgfplots/ymin});
				\coordinate (e2) at (axis cs:638, \pgfkeysvalueof{/pgfplots/ymax});
				
				\draw [dashed] (c1) -- (c2);
				\draw [dashed] (d1) -- (d2);
				\draw [dashed] (e1) -- (e2);
				
				\draw[>=latex, <->] ($(d1) + (rel axis cs:0,0.75)$) -- node [midway,above,align=center,fill=white,yshift=.35cm] {\footnotesize apprentissage\\ \footnotesize\SI{100}{\milli\second}} ($(c1) + (rel axis cs:0,0.75)$);
				\draw[>=latex, <->] ($(c1) + (rel axis cs:0,0.75)$) -- node [midway,above,align=center,yshift=.15cm] {\footnotesize \SI{200}{\milli\second}} ($(e1) + (rel axis cs:0,0.75)$);
				
				%define coordinates at bottom left and top left of rectangle
				\coordinate (r1) at (axis cs:400,0.5);
				\coordinate (r2) at (axis cs:650,3);
				\coordinate (r3) at (axis cs:400,3);
				\coordinate (r4) at (axis cs:650,0.5);			
				\draw (r1) rectangle (r2);
				
				\end{axis}
				%
				\begin{axis}[%
				name=ax2,
				width=.5\linewidth,
				height=.4\linewidth,
				scaled ticks=false,   
				xlabel={temps (\SI{}{\milli\second})},
				axis x line = bottom,
				axis y line = left, 
				xmin=400,xmax=650,
				%ymin=-.05,%ymax=0.25,
				at={($(ax1.south east)+(1cm,0)$)},
				]
				%
				\addplot [
				color=black,
				solid,
				line width=1.pt,
				]
				table [domain=400:650, col sep=comma, x expr=\coordindex, y=actual_trajectory] {misc/test2_burdet/best_candidates.csv};
				%				%
				\addplot [
				color=red,
				dashed,
				line width=1.pt,
				]
				table [domain=400:650, col sep=comma, x expr=\coordindex, y=burdet_candidate] {misc/test2_burdet/best_candidates.csv};
								
				\coordinate (d1) at (axis cs:438, \pgfkeysvalueof{/pgfplots/ymin});
				\coordinate (d2) at (axis cs:438, \pgfkeysvalueof{/pgfplots/ymax});			
				\coordinate (e1) at (axis cs:638, \pgfkeysvalueof{/pgfplots/ymin});
				\coordinate (e2) at (axis cs:638, \pgfkeysvalueof{/pgfplots/ymax});
				
				\draw [dashed] (c1) -- (c2);
				\draw [dashed] (d1) -- (d2);
				\draw [dashed] (e1) -- (e2);
				
				\end{axis}
				% draw dashed lines from rectangle in first axis to corners of second
				\draw [dashed] (r4) -- (ax2.south west);
				\draw [dashed] (r2) -- (ax2.north west);
				
				\end{tikzpicture}
			}
		}		
	\caption{Exemples d'application de la méthodologie de \textcite{burdet_method_2000} pour deux estimations de trajectoire en position sur \SI{200}{\milli\second}. Sur l'ensemble des courbes candidates, seules 45 sont représentées.}
	\label{fig_test_burdet_2000}
\end{figure}
	%%
	%